\documentclass[11pt,a4paper]{article}
\usepackage[utf8]{inputenc}
\usepackage[margin=1in]{geometry}
\usepackage{amsmath,amssymb,amsfonts}
\usepackage{booktabs}
\usepackage{tabularx}
\usepackage{graphicx}
\usepackage{microtype}
\usepackage{hyperref}
\usepackage{cite}
\usepackage{listings}
\usepackage{xcolor}

\hypersetup{
    colorlinks=true,
    linkcolor=blue,
    citecolor=blue,
    urlcolor=teal
}

\lstset{
  basicstyle=\ttfamily\small,
  breaklines=true,
  frame=single,
  backgroundcolor=\color{gray!10}
}

\title{\textbf{Autonomous Forensic Research Auditing Pipeline: Multi-Gate Metascience, 95\% CI Forensics, and Genomic Scrutiny Swarm}}

\author{
  \textbf{Yee How Lau}\thanks{Correspondence to: Yee How Lau, National Heart Centre Singapore, 5 Hospital Drive, Singapore 169609. Email: \texttt{lau.yee.how@nhcs.com.sg}} \\
  \textit{National Heart Centre Singapore} \\
  Singapore
}

\date{\today}

\begin{document}

\maketitle

\begin{abstract}
The replication crisis and the proliferation of paper-mill fabrications, $p$-hacking, and reporting discrepancies pose significant challenges to evidence-based medicine and clinical practice guidelines. Manual forensic auditing of scientific literature is labor-intensive, methodologically fragmented, and difficult to scale across expanding volumes of preprints and published trials. Here, we introduce the \textbf{Autonomous Forensic Research Auditing Pipeline}, an agentic metascience verification swarm integrated with the Model Context Protocol (MCP) and large language model architectures. The framework enforces a strict non-compensatory two-tier triage protocol: \textbf{Tier 1} applies a 4-gate lexicographic hurdle covering provenance and retraction status, tortured phrases and wet-lab reagent validation, universal table granularity (GRIM, GRIMMER, DEBIT, SPRITE, and Benford's Law), 95\% confidence interval consistency and variance compression forensics, and multi-omics genomic grounding (linkage disequilibrium support towers and volcano plot FDR control). Studies clearing Tier 1 without triggering fatal mathematical impossibility kill-switches proceed to \textbf{Tier 2}, where empirical evidence across risk-of-bias domains, open code/data artifacts, and statistical precision is synthesized into an objective Bayesian Posterior Credibility Probability via empirical likelihood ratio multipliers ($LR_i$). To democratize reproducibility and accelerate peer scrutiny, the system automatically generates structured consensus reports, neutral PubPeer-compliant review briefs, and Twitter/X dissemination threads. The complete framework is open-source, versioned on Zenodo (\texttt{doi:10.5281/zenodo.23861362}), and accessible on GitHub.
\end{abstract}

\vspace{1em}
\noindent \textbf{Keywords:} Metascience, Forensic Audit, Statistical Fabrication, 95\% Confidence Intervals, Multi-Omics Verification, Bayesian Likelihood Ratio, PubPeer.

\section{Introduction}
Scientific progress depends on the premise that published data, statistical derivations, and methodological descriptions are internally valid and mathematically plausible. However, empirical audits demonstrate that reporting discrepancies, phantom citations, tortured phrases generated by spinning software, and outright data fabrication are widespread across biomedical, psychological, and clinical disciplines~\cite{cabanac2021tortured, nuijten2016statcheck, carlisle2017data}.

Traditional peer review remains largely qualitative and ill-equipped to systematically verify numeric granularity, such as checking whether reported integer means align with denominator arithmetic (the GRIM test~\cite{brown2017grim}), whether reported standard deviations are mathematically permissible (the GRIMMER test~\cite{allard2018grimmer}), whether binary proportions match reported sample sizes (the DEBIT test~\cite{heathers2019debit}), whether test statistics match reported $p$-values (\texttt{statcheck}~\cite{nuijten2016statcheck}), or whether claimed 95\% confidence intervals (CIs) exhibit variance compression or mathematical asymmetry~\cite{altman2011confidence, guyatt2011grade}.

To address these vulnerabilities, we developed the \textit{Autonomous Forensic Research Auditing Pipeline}, an open-source, deterministic agentic framework engineered to execute comprehensive, multi-modal audits of published papers, preprints, and clinical registries. The pipeline synthesizes algorithmic error detection with an overarching Bayesian credibility scoring framework.

\section{System Architecture and Methodology}
The auditing engine operates under a non-compensatory two-tier model: no study can compensate for a fatal mathematical or biological impossibility through prestige, high journal impact factors, or positive clinical endpoints.

\subsection{Tier 1: Multi-Gate Lexicographic Triage}
\begin{enumerate}
    \item \textbf{Gate 1: Provenance, Retraction, and Registry Verification.}
    Queries Crossref, OpenAlex, PubMed, and Retraction Watch database indices via MCP servers to detect formal expressions of concern, errata, or retractions. Confirms pre-registration consistency on ClinicalTrials.gov or ISRCTN registries.
    
    \item \textbf{Gate 2: Tortured Phrases and Wet-Lab Reagent Forensics.}
    Scans the corpus against the Problematic Paper Screener corpus of tortured phrases~\cite{cabanac2021tortured} (e.g., ``counterfeit consciousness'' for artificial intelligence, ``colossal bowel'' for large intestine). Cross-references cell lines against the International Cell Line Authentication Committee (ICLAC) database of cross-contaminated or misidentified cell lines and validates nucleotide primer sequences against BLAST/NCBI nucleotide registries.
    
    \item \textbf{Gate 3: Universal Table, Granularity, and 95\% CI Forensics.}
    \begin{itemize}
        \item \textbf{Granularity Auditing:} Executes the GRIM~\cite{brown2017grim}, GRIMMER~\cite{allard2018grimmer}, and DEBIT~\cite{heathers2019debit} tests on all reported summary tables. Applies SPRITE~\cite{heathers2018sprite} to evaluate bounded distributions and tests first/second-digit distributions against Benford's Law~\cite{benford1938law}.
        \item \textbf{95\% Confidence Interval Forensics:} For all reported effect sizes (mean differences, hazard ratios, odds ratios), calculates reconstructed standard errors ($SE = (UL - LL) / 3.9199$)~\cite{altman2011confidence} and validates symmetry on linear or log-scales. Reconstructs test statistics $Z = \text{Estimate} / SE$ and computes recalculated two-sided $p$-values ($p_{\text{recalc}}$).
        \item \textbf{Variance Compression vs. Imprecision:} Automatically identifies implausibly narrow CIs ($\text{implied } SD = SE \times \sqrt{N} < 25\%$ of reference clinical variance, indicating standard error misreported as confidence interval) or excessively wide CIs spanning minimal clinically important differences (MCID) and failing Cochrane GRADE optimal information size thresholds~\cite{guyatt2011grade}.
    \end{itemize}
    
    \item \textbf{Gate 3.5: Multi-Omics and Chart Density Grounding.}
    Audits clinical scatterplots for density over-extrapolation ($<10\%$ of observations in headline therapeutic ranges). In genome-wide association studies (GWAS), evaluates Manhattan plots for linkage disequilibrium (LD) support towers ($\ge 2$ flanking variants within $\pm 500$ kb with $p < 10^{-4}$) to eliminate isolated single-SNP chimneys. In transcriptomics, verifies that volcano plots use Benjamini-Hochberg false discovery rate (FDR) control rather than unadjusted raw $p$-values.
\end{enumerate}

\subsection{Kill-Switch Rule}
If a study exhibits any fatal violation—such as a \texttt{statcheck} decision error ($p_{\text{reported}} < 0.05$ while $p_{\text{recalculated}} \ge 0.05$), $\ge 2$ GRIM/DEBIT impossible means, non-existent dbSNP genetic identifiers, or verified retraction—the pipeline halts execution immediately. A \texttt{HARD REJECT (0\% Credibility)} consensus is issued without proceeding to Bayesian scoring.

\subsection{Tier 2: Bayesian Likelihood Ratio Credibility Score}
Papers clearing Tier 1 are evaluated through a formal Bayesian updating model. Starting from a skeptical baseline prior probability $P_0 = 0.30$ ($\text{Prior Odds} = 0.30 / (1 - 0.30) = 0.4286$), empirical evidence multipliers ($LR_i$) are iteratively applied:
\begin{equation}
    \text{Posterior Odds} = \text{Prior Odds} \times \prod_{i=1}^{k} LR_i
\end{equation}
\begin{equation}
    P(\text{Credible} \mid \text{Evidence}) = \frac{\text{Posterior Odds}}{1 + \text{Posterior Odds}}
\end{equation}

Representative empirical likelihood ratios include:
\begin{itemize}
    \item Open raw data and reproducible analysis code: $LR = 2.00$
    \item Independent pre-registered replication: $LR = 4.00$
    \item Primary $p$-value in the $0.041$--$0.049$ $p$-hacking margin: $LR = 0.25$
    \item PROBAST~\cite{wolff2019probast} or Cochrane RoB 2~\cite{sterne2019rob2} ``Some Concerns'' rating: $LR = 0.50$
    \item Cochrane RoB 2 ``High Risk of Bias'': $LR = 0.15$
    \item Implausibly narrow 95\% CI ($SE$-as-CI reporting): $LR = 0.15$
    \item Severe 95\% CI imprecision spanning clinical harm and benefit~\cite{guyatt2011grade}: $LR = 0.40$
    \item Unreconciled sample size transitions ($N$ drop across tables): $LR = 0.30$
\end{itemize}

Posterior probabilities are mapped deterministically to four standardized certainty tiers:
\begin{itemize}
    \item \textbf{HIGH INTEGRITY}: $P \ge 85\%$ (GRADE High Certainty)
    \item \textbf{QUALIFIED PASS}: $45\% \le P < 85\%$ (GRADE Moderate Certainty)
    \item \textbf{SKEPTICAL CAUTION}: $15\% \le P < 45\%$ (GRADE Low Certainty)
    \item \textbf{HIGH RISK / CRITICAL CONCERN}: $P < 15\%$ (GRADE Very Low Certainty)
\end{itemize}

\section{Automated Deliverable Generation}
To maximize real-world transparency and facilitate post-publication peer review, every completed audit generates three persistent artifacts:
\begin{enumerate}
    \item \textbf{Consensus Audit Report (\texttt{audit\_[ID].md}):} A five-section technical dossier detailing provenance, granularity checks, 95\% CI forensic derivations, risk of bias assessments, and the full Bayesian likelihood ratio chain.
    \item \textbf{PubPeer Review Brief (\texttt{pubpeer\_[ID].md}):} A neutral, strictly evidence-based comment formatted for direct submission to the PubPeer community platform.
    \item \textbf{Social Dissemination Thread (\texttt{tweet\_[ID].md}):} A four-post structured Twitter/X disclosure thread tagged for scientific integrity investigators and community sleuths.
\end{enumerate}

\section{Code and Software Availability}
The pipeline source code, R validation scripts, MCP tool servers, and prompt specifications are open-source and freely accessible under the MIT License:
\begin{itemize}
    \item \textbf{Source Code Repository:} \url{https://github.com/lauyeehow1986-hub/forensic_audit}
    \item \textbf{Archival Release (Zenodo DOI):} \href{https://doi.org/10.5281/zenodo.23861362}{10.5281/zenodo.23861362}
\end{itemize}

\bibliographystyle{unsrt}
\bibliography{references}

\end{document}
